\chapter{7B 推理计算账本：FLOPs、字节、KV cache 与延迟上限}

\section{为什么必须算账}

讨论本地 7B 推理时，最常见的空话是“优化矩阵乘”“减少显存”“提高 tokens/s”。这些话如果不落到 FLOPs、字节、cache、带宽和延迟，就无法指导实现。本章用一个典型 7B decoder-only GQA 模型建立账本。具体模型可能略有差异，但方法通用。

\begin{lstlisting}[numbers=none,caption={典型 7B 账本参数}]
layers L = 32
hidden H = 4096
query_heads n_q = 32
kv_heads n_kv = 8
head_dim d = 128
intermediate I = 11008
vocab V = 32000
context S = 4096
\end{lstlisting}

这些数字不是装饰。它们决定每个 token 要读多少权重、KV cache 占多少内存、attention 随上下文怎样增长、CPU/GPU 后端为什么在 prefill 和 decode 阶段表现不同。

\section{参数量账本}

先算一层主要权重。Q projection 输出 \code{n\_q*d=4096}，K/V projection 各输出 \code{n\_kv*d=1024}。

\begin{center}
\begin{tabular}{p{0.26\linewidth}p{0.30\linewidth}p{0.28\linewidth}}
\toprule
权重 & 形状 & 参数量 \\
\midrule
\code{Wq} & \code{4096 x 4096} & \code{16.78M} \\
\code{Wk} & \code{1024 x 4096} & \code{4.19M} \\
\code{Wv} & \code{1024 x 4096} & \code{4.19M} \\
\code{Wo} & \code{4096 x 4096} & \code{16.78M} \\
\code{Wgate} & \code{11008 x 4096} & \code{45.09M} \\
\code{Wup} & \code{11008 x 4096} & \code{45.09M} \\
\code{Wdown} & \code{4096 x 11008} & \code{45.09M} \\
\bottomrule
\end{tabular}
\end{center}

一层线性权重大约：

\[
16.78 + 4.19 + 4.19 + 16.78 + 45.09 + 45.09 + 45.09 \approx 177.2M
\]

32 层约 \code{5.67B} 参数。加上 embedding、lm head、norm 和少量元数据，就接近 7B 级别。这个账本直接解释了 decode 为什么常常被权重读流支配：每生成一个 token，主要线性权重基本都要参与一次。

\section{权重字节账本}

参数量乘以每个权重的字节数，就是权重读流的粗略下界。

\begin{center}
\begin{tabular}{p{0.18\linewidth}p{0.28\linewidth}p{0.34\linewidth}}
\toprule
格式 & 7B 权重字节 & 备注 \\
\midrule
fp16 & 约 \code{14 GB} & 精度高，带宽压力大 \\
int8 & 约 \code{7 GB} & 还要加 scale/zero point \\
int4 & 约 \code{3.5 GB} & metadata 和 packing 会增加一些 \\
\bottomrule
\end{tabular}
\end{center}

若每个 decode token 都要流过大部分权重，那么 CPU 上限首先受内存带宽约束。假设有效内存带宽是 \code{100 GB/s}，int4 权重加 metadata 后每 token 读 \code{3.7 GB}，只看权重的理论上限也只有：

\[
\frac{100}{3.7} \approx 27 \text{ tokens/s}
\]

这还没有算 KV cache 读取、activation、线程同步、解包、scale、sampler、allocator 和调度开销。真实速度低于这个上限并不奇怪。这个公式比单纯说“CPU 慢”更有用，因为它告诉我们应该优先减少权重字节、改善 packing、提高带宽利用率，而不是盲目堆线程。

\section{decode 每层 FLOPs 账本}

一次 decode 只处理一个新 token。线性层 MAC 数大约等于权重参数量。每层主要线性 MAC：

\[
QKV = H(H+2n_{kv}d)=4096(4096+2048)\approx25.17M
\]

\[
O = H^2 \approx16.78M
\]

\[
MLP = 3HI = 3 \times 4096 \times 11008 \approx135.27M
\]

总计每层约：

\[
25.17 + 16.78 + 135.27 \approx 177.22M \text{ MACs}
\]

若把一次乘加按 2 FLOPs 计，每层约 \code{354 MFLOPs}，32 层约 \code{11.3 GFLOPs/token}。attention 的点积和加权求和还要另算，随上下文长度 \code{S} 增长：

\[
attention\_macs \approx 2 \times n_q \times S \times d
\]

\code{S=4096} 时，一层 attention 约：

\[
2 \times 32 \times 4096 \times 128 \approx 33.55M \text{ MACs}
\]

32 层约 \code{1.07B MACs}。所以在长上下文下，attention 计算也很可观；但在很多 CPU decode 场景，低比特权重的解包和权重/KV 字节流仍然是更早撞到的墙。

\section{KV cache 容量账本}

KV cache 每个 token、每层要保存 K 和 V。GQA 下保存的是 \code{n\_kv} 个 head，不是 \code{n\_q} 个 head。

\[
kv\_bytes\_per\_token\_per\_layer =
2 \times n_{kv} \times d \times bytes
\]

fp16 下每个元素 2 字节：

\[
2 \times 8 \times 128 \times 2 = 4096 \text{ bytes}
\]

32 层就是每个 token \code{128 KiB}。上下文 \code{4096} 时，一个 session 的 KV cache 是：

\[
4096 \times 128KiB = 512MiB
\]

这个数字对业务开发非常关键。一个本地服务如果同时接 8 个 \code{4096} context 的会话，仅 KV cache 就可能接近 \code{4 GiB}，还不包括权重、packed weight、activation arena、workspace、线程栈和系统余量。

\begin{center}
\begin{tabular}{p{0.18\linewidth}p{0.28\linewidth}p{0.30\linewidth}}
\toprule
context & 单 session fp16 KV & 8 session fp16 KV \\
\midrule
\code{1024} & \code{128 MiB} & \code{1 GiB} \\
\code{2048} & \code{256 MiB} & \code{2 GiB} \\
\code{4096} & \code{512 MiB} & \code{4 GiB} \\
\code{8192} & \code{1 GiB} & \code{8 GiB} \\
\bottomrule
\end{tabular}
\end{center}

这就是为什么 admission control 必须在请求进入前估算 KV cache，而不是等到 decode 过程中分配失败。

\section{KV cache 读取账本}

容量只是第一半。decode 每步还要读取历史 K/V。当前 token 在每层 attention 中读取长度为 \code{S} 的 K 和 V：

\[
kv\_read\_bytes\_per\_layer =
S \times 2 \times n_{kv} \times d \times bytes
\]

\code{S=4096}、fp16 时每层约 \code{16 MiB}，32 层约 \code{512 MiB/token}。也就是说，在长上下文下，即使用 int4 权重，每生成一个 token 还要额外读约半 GB 的 KV cache。这解释了两个现象：

\begin{itemize}
  \item 上下文越长，decode token 越慢，不只是 attention FLOPs 增加，KV 读字节也线性增加。
  \item KV cache 量化可能有效，因为它减少的是每步反复读取的历史状态，而不是只减少静态权重。
\end{itemize}

但 KV 量化比权重量化更敏感。权重是固定的，可以离线校准；KV 是动态 activation，每个请求、每个位置都不同。若 scale 选择不好，attention score 会被放大后进入 softmax，少量误差可能改变 top-k。后面的低比特 KV 章节会专门展开。

\section{prefill 和 decode 不是同一账本}

prefill 处理整个 prompt。若 prompt 长度 \code{T=512}，一层 linear 变成 \code{[512,H] x [H,O]} 的矩阵乘。权重被同一批 token 复用，GEMM 能更好利用计算单元。decode 每步 \code{T=1}，权重复用少，容易被权重读流限制。

\begin{center}
\begin{tabular}{p{0.18\linewidth}p{0.34\linewidth}p{0.32\linewidth}}
\toprule
阶段 & 主要形态 & 常见瓶颈 \\
\midrule
prefill & 大 GEMM、批量 attention & 算力、显存带宽、attention tile \\
decode & GEMV/小 GEMM、KV 读取 & 权重带宽、KV 带宽、launch/同步 \\
\bottomrule
\end{tabular}
\end{center}

这就是为什么一个后端可能 prefill 很快、decode 一般；也可能 CPU decode 经过 int4 和线程优化后可用，但长 prompt prefill 明显慢。报告必须分开写 TTFT、prefill tokens/s 和 decode tokens/s。

\section{roofline：用带宽和算力约束 tokens/s}

粗略上限可以用 roofline 思路判断。对每个 decode token，假设需要读：

\begin{lstlisting}[numbers=none,caption={decode token 的粗略字节项}]
weight_bytes_per_token
kv_read_bytes_per_token(context)
activation_bytes
metadata_bytes
temporary_workspace_bytes
\end{lstlisting}

带宽上限：

\[
tokens/s \le
\frac{effective\_bandwidth}{bytes\_per\_token}
\]

算力上限：

\[
tokens/s \le
\frac{effective\_flops}{flops\_per\_token}
\]

实际上限取二者较小者，再扣除线程调度、解包、sampler、同步和 cache miss。举例：int4 权重 \code{3.7 GB/token}，长上下文 KV \code{0.5 GB/token}，总计约 \code{4.2 GB/token}。若有效带宽 \code{100 GB/s}：

\[
\frac{100}{4.2}\approx23.8 \text{ tokens/s}
\]

若同一机器由于解包、NUMA、TLB、线程调度只能达到 \code{55 GB/s} 有效带宽，上限就是 \code{13 tokens/s}。这比“开更多线程”更能解释问题：线程只能提高并行利用率，不能凭空减少每 token 字节数。

\section{量化为什么有效，为什么也危险}

权重量化主要减少静态权重字节。int8 把 fp16 权重从 2 字节变成 1 字节；int4 变成 0.5 字节，再加 scale metadata。若 decode 被权重带宽限制，低比特能直接提高上限。

但量化不是免费午餐。计算中会出现：

\begin{itemize}
  \item 解包成本：int4 需要从 byte 中取 nibble。
  \item scale 成本：per-group 或 per-channel scale 要参与反量化。
  \item accumulator 精度：低比特乘法需要更宽累加或合理 requantize。
  \item 误差传播：Q/K/V、MLP、lm head 对最终 logits 的敏感度不同。
  \item metadata 读流：scale、zero point 和 group table 也占带宽。
\end{itemize}

所以量化报告不能只写“模型从 14GB 变成 3.7GB”。它至少要写：

\begin{lstlisting}[numbers=none,caption={量化收益报告字段}]
original_weight_bytes
packed_weight_bytes
scale_metadata_bytes
decode_tokens_per_s
prefill_tokens_per_s
operator_error
logits_topk_overlap
generation_diff_under_fixed_seed
\end{lstlisting}

低比特格式能不能进入业务服务，必须由容量、速度和质量三者共同决定。

\section{并发容量账本}

业务服务要算的是总内存，不是单请求。一个粗略公式：

\begin{lstlisting}[numbers=none,caption={服务内存预算}]
total =
  model_weight_bytes
  + packed_weight_bytes
  + backend_static_workspace
  + concurrent_sessions * kv_cache_bytes_per_session
  + concurrent_sessions * activation_arena_bytes
  + trace_buffers
  + safety_margin
\end{lstlisting}

假设 int4 packed 权重和 metadata 约 \code{4.0 GB}，静态 workspace \code{0.5 GB}，每个 \code{4096} context session 的 KV 为 \code{512 MiB}，每个 session arena 和 trace 估计 \code{128 MiB}，机器可给进程使用 \code{12 GB}。最多并发不能简单算 \code{12/0.5=24}，而应先扣静态项：

\[
12 - 4.0 - 0.5 - safety \approx 6.5GB
\]

每 session 约 \code{640 MiB}，最多约 10 个，还要考虑 allocator 碎片和峰值。因此 admission control 应该保守，并按 \code{max\_context\_tokens} 预留，而不是按当前 prompt 长度乐观估计。

\section{把账本写进 trace}

计算账本如果只在书里，工程上仍然不可用。trace event 应该包含估计字节和实际耗时：

\begin{lstlisting}[numbers=none,caption={账本相关 trace 字段}]
TraceEvent:
  stage
  layer_id
  segment_id
  context_length
  estimated_flops
  estimated_weight_bytes
  estimated_kv_read_bytes
  estimated_kv_write_bytes
  elapsed_ns
  backend
  quant_profile_id
\end{lstlisting}

有了这些字段，报告就能回答：

\begin{itemize}
  \item decode 变慢是上下文变长导致 KV 读增加，还是后端 fallback。
  \item int4 变快是否符合权重字节下降的预期。
  \item prefill 慢是 GEMM 吞吐问题，还是 tokenizer/prepare 被算进了 TTFT。
  \item 某层异常慢是否因为 layout 转换或 workspace 分配。
\end{itemize}

\section{本章验收线}

读完本章，应该能自己给一个 7B 模型算粗账：

\begin{itemize}
  \item 根据 \code{L,H,n\_q,n\_kv,d,I,S} 估算每层参数量和 decode FLOPs。
  \item 估算 fp16、int8、int4 权重字节，并用带宽上限推 tokens/s。
  \item 计算单 session KV cache 容量和多 session 并发内存。
  \item 区分 prefill 的 GEMM 账本和 decode 的 GEMV/KV 读账本。
  \item 解释为什么长上下文会增加 KV 读字节和 attention 计算。
  \item 把 FLOPs、权重字节、KV 字节和耗时写进 trace，而不是只报最终速度。
\end{itemize}
